\chapter{Örnekleme Yöntemleri, Temsil ve Yanlılık}
\label{ch:sampling-design}

\begin{prismbox}[PrismLight]{Öğrenme çıktıları}
Bu bölümün sonunda:
\begin{itemize}
  \item hedef evren, örnekleme çerçevesi ve örneklem arasındaki farkı açıklayabilecek,
  \item araştırma bağlamına uygun örnekleme yöntemini seçebilecek,
  \item seçilme olasılığı, örnekleme aralığı ve örnekleme ağırlığını hesaplayabilecek,
  \item kümeleşme, yanıtlamama ve kapsam hatasının sonuçlarını değerlendirebileceksiniz.
\end{itemize}
\end{prismbox}

\begin{prerequisitebox}
Hedef evren, erişilebilir evren, gözlem birimi ve örneklem kavramları
bilinmelidir. Örnekleme değişkenliği ile sistematik yanlılık arasındaki fark
hatırlanmalıdır.
\end{prerequisitebox}

Hedef evren ve kapsam kavramları Bölüm~\ref{sec:population-target} ile
Bölüm~\ref{sec:population-coverage}'de; örnekleme değişkenliği ile yanlılık
ayrımı Bölüm~\ref{sec:sampling-bias-variance}'ta hazırlanmıştır.

\section{Örnekleme çerçevesi ve kapsam hatası}
\label{sec:sampling-frame}

Araştırmacının hakkında sonuç üretmek istediği topluluk \textbf{hedef evren},
seçim yapılırken fiilen erişilebilen birimlerin listesi \textbf{örnekleme
çerçevesi}, ölçülen birimler ise \textbf{örneklem}dir. Bu üç küme aynı değilse
genelleme yapılırken kapsam hatası oluşabilir.

Örneğin bir üniversitedeki bütün öğrenciler hedeflenirken anket bağlantısının
yalnız belirli bir dersin çevrim içi sayfasında paylaşılması, örnekleme
çerçevesini o dersi alan öğrencilerle sınırlar.

\section{Olasılıklı örnekleme yöntemleri}
\label{sec:sampling-probability}

\begin{description}[style=nextline,leftmargin=3.7cm,labelwidth=3.4cm]
  \item[Basit rastgele örnekleme] Her birim bilinen ve eşit seçilme şansına sahiptir.
  \item[Sistematik örnekleme] Rastgele başlangıçtan sonra listedeki her $k$'ıncı birim seçilir.
  \item[Tabakalı örnekleme] Evren önemli alt gruplara ayrılır ve her tabakadan örnek alınır.
  \item[Küme örneklemesi] Bireyler yerine sınıf, okul veya kurum gibi doğal kümeler seçilir.
  \item[Çok aşamalı örnekleme] Önce kümeler, sonra kümeler içindeki alt birimler seçilir.
\end{description}

Tabakalama, alt grupların temsilini güvenceye almak için; küme örneklemesi ise
dağınık bir evrende veri toplama maliyetini azaltmak için kullanılabilir.
Küme içindeki bireyler birbirine benzediğinde etkili bilgi miktarı, aynı
sayıdaki bağımsız bireyden daha düşük olur. Bu tasarımların seçim olasılıkları,
ağırlıklar ve varyans üzerindeki etkileri \citet{cochran1977} ve
\citet{lohr2021} tarafından ayrıntılı biçimde ele alınır.

Şekil~\ref{fig:strata-clusters} iki yöntemin seçim birimini karşılaştırır.
Solda her tabakadan ikişer kişi; sağda ise bir küme seçilerek o kümedeki
altı kişinin tamamı alınmıştır. Her iki örneklem altı kişilik olsa da
seçim mekanizmaları farklıdır. Sağdaki çizim tek aşamalı küme örneklemesidir;
çok aşamalı tasarımda seçilen kümenin içinde ayrıca birey seçimi yapılabilir.

\begin{figure}[H]
\centering
\begin{tikzpicture}[font=\footnotesize]
  \foreach \shift/\heading in {-2.7/Tabakalı örnekleme,2.7/Küme örneklemesi}{
    \begin{scope}[xshift=\shift cm]
      \node[font=\small\bfseries] at (0,0.6) {\heading};
      \foreach \row in {1,2,3}{
        \draw[PrismBlue!65,rounded corners=1mm]
          (-2.35,{-0.85*\row-0.27}) rectangle (2.35,{-0.85*\row+0.27});
        \foreach \col in {1,...,6}
          \draw[PrismBlue,fill=white]
            ({-0.75+0.48*(\col-1)},{-0.85*\row}) circle (2.8pt);
      }
    \end{scope}
  }
  \foreach \row in {1,2,3}{
    \node[anchor=west] at (-4.9,{-0.85*\row}) {Tabaka \row};
    \node[anchor=west] at (0.5,{-0.85*\row}) {Küme \row};
    \foreach \col in {2,5}
      \fill[PrismBlue] ({-3.45+0.48*(\col-1)},{-0.85*\row}) circle (2.8pt);
  }
  \foreach \col in {1,...,6}
    \fill[PrismBlue] ({1.95+0.48*(\col-1)},-1.7) circle (2.8pt);
  \node[align=center] at (-2.7,-3.25) {Her tabakadan birey seçimi};
  \node[align=center] at (2.7,-3.25) {Seçilen kümenin tüm bireyleri};
  \fill[PrismBlue] (-2.8,-3.95) circle (2.8pt);
  \node[anchor=west] at (-2.6,-3.95) {Seçilen};
  \draw[PrismBlue,fill=white] (0.3,-3.95) circle (2.8pt);
  \node[anchor=west] at (0.5,-3.95) {Seçilmeyen};
\end{tikzpicture}
\caption{Aynı örneklem hacmi, farklı seçim mekanizması: tabakalı ve tek
aşamalı küme örneklemesi. Dolu noktalar seçilen bireylerdir.}
\label{fig:strata-clusters}
\end{figure}

\section{Olasılıksız örnekleme}
\label{sec:sampling-nonprobability}

Kolayda, gönüllü ve amaçlı örnekleme bazı araştırma bağlamlarında gerekli
olabilir; ancak birimlerin seçilme olasılıkları bilinmediği için örnekleme
hatası klasik formüllerle tam olarak hesaplanamaz. Sonuçların hangi gruba
genellenebileceği açıkça sınırlandırılmalıdır.

\section{Yanlılık kaynakları}
\label{sec:sampling-bias-sources}

\begin{itemize}
  \item \textbf{Kapsam yanlılığı:} Çerçeve hedef evrenin bazı kesimlerini dışlar.
  \item \textbf{Yanıtlamama yanlılığı:} Yanıt verenler vermeyenlerden sistematik olarak farklıdır.
  \item \textbf{Gönüllülük yanlılığı:} Güçlü görüşü olan kişiler katılmaya daha isteklidir.
  \item \textbf{Ölçüm yanlılığı:} Soru biçimi, araç veya uygulayıcı yanıtları sistematik etkiler.
  \item \textbf{Seçim sonrası dışlama:} Analiz ölçütleri belirli sonuçları aşırı temsil eder.
\end{itemize}

\begin{mistakebox}
Rastgele örnekleme ile rastgele atama aynı işlem değildir. Rastgele örnekleme
evrene genellemeyi, rastgele atama ise gruplar arasında nedensel karşılaştırmayı
destekler.
\end{mistakebox}

\section{Örneklem hacmi ve temsil}

Büyük örneklem rastgele örnekleme hatasını azaltır; fakat sistematik yanlılığı
otomatik olarak gidermez. Örneklem hacmi belirlenirken hedeflenen hassasiyet,
değişkenlik, güven düzeyi, araştırma tasarımı ve ulaşılabilir kaynaklar birlikte
değerlendirilir.

Basit rastgele örneklemede ortalamanın standart hatası yaklaşık
\[
\SE(\bar X)=\frac{s}{\sqrt n}
\]
olduğundan belirsizliği yarıya indirmek yaklaşık dört kat gözlem gerektirir.

\section{Seçilme olasılığı ve örnekleme ağırlığı}
\label{sec:sampling-weights}

Basit rastgele örneklemede $N$ birimlik bir evrenden $n$ birim seçiliyorsa her
birimin örnekleme dâhil edilme olasılığı
\[
\pi_i=\frac{n}{N}
\]
olur. Bu olasılığın tersi $w_i=1/\pi_i$ temel örnekleme ağırlığıdır. Örneğin
2.000 öğrenciden 200'ü eşit olasılıkla seçiliyorsa $\pi_i=0{,}10$ ve her
öğrencinin temel ağırlığı 10'dur; örnekteki bir öğrenci evrende yaklaşık on
öğrenciyi temsil eder.

Tabakalar farklı oranlarda örneklenirse seçilme olasılıkları ve ağırlıklar da
tabakaya göre değişir. Ağırlıkları göz ardı etmek, fazla örneklenen küçük
grupların evren özetleri üzerindeki etkisini gereğinden fazla artırabilir.

\section{Kümeleşme ve tasarım etkisi}
\label{sec:sampling-deff}

Küme örneklemesinde aynı sınıf veya okul içindeki bireyler birbirine
benzeyebilir. Ortalama küme büyüklüğü $m$ ve küme içi korelasyon $\rho$ için
basit bir tasarım etkisi yaklaşımı
\[
\DEFF\approx1+(m-1)\rho
\]
biçimindedir. Tasarım etkisinin 1'den büyük olması, aynı sayıda tamamen
bağımsız gözleme göre varyansın arttığını gösterir. Yaklaşık etkin örneklem
hacmi $n_{\mathrm{etkin}}\approx n/\DEFF$ ile yorumlanabilir.

Bu formül öğretici bir yaklaşımdır; gerçek analizde eşit olmayan küme
büyüklükleri, tabakalama, ağırlıklandırma ve çok aşamalı seçim birlikte
değerlendirilmelidir.

\section{Tamamlayıcı vaka: örneklem dosyası çerçeve değildir}
\label{sec:v2-nhanes}

\textbf{Soru ve hedef.} NHANES 2017--2018 araştırmasında ABD'nin 50
eyaleti ve Columbia Bölgesi'nde yaşayan, kurumlarda yaşamayan sivil
nüfus hedeflenir. Bu vaka için sorumuz daha dardır: bu nüfusun 20 yaş
ve üzerindeki kesiminde, 60 yaş ve üzerindekilerin oranı nedir?
Dönem ve yaş sınırları analiz tanımının parçasıdır; hedef bugünkü
nüfus veya Türkiye nüfusu değildir. Bu bir müdahale etkisi sorusu da
değildir. Veri kaynağı demografi dosyası \texttt{DEMO\_J}'dir
\citep{nchs2020demoj,nchssampledesign}.

\textbf{Çerçeveden kayda.} Seçim, coğrafi birincil örnekleme
birimlerinden segmentlere, hanelere ve kişilere ilerleyen çok aşamalı
bir düzene dayanır. Hane listeleri ve tarama işlemleri, kişilerin
seçilebilmesini sağlar. Kamuya açık yanıtlayıcı dosyası bu çerçevelerin
yerine geçmez: seçilmeyen veya yanıt vermeyen bütün birimleri içermez
\citep{nchssampledesign}. B06'daki yapay çerçeveden bilgisayarla seçim
yapmak, burada uygulanmış saha tasarımını yeniden gerçekleştirmek değildir.

\textbf{Analiz tanımı.} Gözlem birimi kişidir. Yaş için
\texttt{RIDAGEYR}, görüşme ağırlığı için \texttt{WTINT2YR}; varyans
hesabı için \texttt{SDMVSTRA} ve \texttt{SDMVPSU} kullanılacaktır.
Son iki değişken gizliliği koruyan varyans birimleridir; gerçek coğrafi
yer adları veya doğrudan saha çerçevesi olarak okunmaz. Yaştaki 80
kodu 80 ve üzerini birleştirir; seçtiğimiz 60 yaş eşiğini belirsiz
hale getirmez \citep{nchs2020demoj}.

\begin{definitionbox}
\textbf{İki farklı özet.} $d_i=\Ind{\text{yaş}_i\geq20}$ ve
$y_i=\Ind{\text{yaş}_i\geq60}$ olsun. $w_i$ görüşme ağırlığıdır.
\[
\widehat p_w=\frac{\sum_i w_i d_i y_i}{\sum_i w_i d_i},
\qquad
\widehat p_0=\frac{\sum_i d_i y_i}{\sum_i d_i}.
\]
İlk ifade hedef alt nüfusun oranını tahmin etmek için kullanılacak
ağırlıklı orandır. İkincisi dosyadaki yetişkin yanıtlayıcıların
bileşimini betimler. İki sonuç karşılaştırılabilir; fakat tercih edilen
sayısal sonucu verdiği için ağırlık seçilmez veya kaldırılmaz.
\end{definitionbox}

\textbf{Ağırlık neyi düzeltir?} Nihai ağırlık, temel seçim ağırlığının
yanı sıra yanıtsızlık ve nüfus toplamlarına uyarlama işlemlerini de
içerir. Bu nedenle $1/w_i$ değerini özgün seçilme olasılığı diye
raporlamayız. Yalnız görüşme değişkenlerini kullanan bu soru için
muayene ağırlığına geçmeyiz. Ağırlıklandırma, çerçeve kapsamı veya
yanıtsızlıktan doğan her yanlılığın giderildiğini garanti etmez
\citep{nchsweighting}.

\textbf{Belirsizlik nasıl hesaplanacak?} Tabaka, küme ve ağırlıkları
birlikte kullanan Taylor doğrusallaştırmasıyla tasarıma uygun standart
hata hesaplanmalıdır. Yetişkinler bir \emph{alt evren} olarak tanımlanır;
tasarım kurulmadan çocuk kayıtları silinmez. Yalnız ağırlıklı oranı
bulmak veya $\sqrt{\widehat p(1-\widehat p)/n}$ kullanmak bu hesabın
yerine geçmez \citep{nchsvariance}. Tahmin ile güven aralığı, farklı
tamamlanma kontrolleridir.

\begin{researchbox}
\textbf{Vakanın mevcut kapsamı.} Kaynak belgelerine dayalı tasarım
okumasını, görevlerin ardından gelen sayısal uygulama tamamlar.
Dosya kimliği, yaş kodları, pozitif ağırlıklar, tabaka--PSU yapısı ve
seçili alanlarda eksiksizlik kontrol edilmiştir. Diğer değişkenlerdeki
yapısal eksiklikler yüzünden bütün sütunlarda topluca tam kayıt
seçilmemiştir. Beklenmedik girdide kod otomatik dışlama yerine durur.
Bu dar yaş oranı uygulaması NHANES'in bütün analizleri için kabul
veya yayın onayı değildir.
\end{researchbox}

\subsection{Karar görevleri ve gerekçeli kontrol}
Önce aşağıdaki beş görevi yanıtlayın; ardından kontrol yanıtlarıyla
karşılaştırın. Sayısal veri dosyası olmadan da tasarım kararları
değerlendirilebilir.

\begin{enumerate}
\item Hedef nüfusu, alt evreni ve dönemi yazın. Kurumlarda yaşayanları
aynı hedefe dahil etmek neden yalnız ağırlık değişikliği değildir?
\item Bir öğrenci ``\texttt{DEMO\_J}'den rastgele 200 satır seçersem
özgün örnekleme çerçevesini elde ederim'' diyor. Yanılgıyı açıklayın.
\item İki oran formülünün paydalarını yorumlayın. Nihai ağırlığın tersi
neden tek başına saha seçiminin olasılığını vermez?
\item ``Çocukları siler, ağırlıklı oranı bulur, bağımsız gözlem
formülüyle standart hata hesaplarım'' planını düzeltin.
\item Sonuç raporu için iki kapsam cümlesi yazın: biri kime ve hangi
döneme ilişkin tahmin yapılacağını, diğeri hangi iddianın
kurulamayacağını belirtsin. Henüz hesaplanmamış değer yazmayın.
\end{enumerate}

\textbf{Kontrol yanıtları.}
\begin{enumerate}
\item Araştırmanın hedefindeki 20 yaş ve üzeri kişiler alt evrendir;
dönem 2017--2018'dir. Kurum nüfusunu eklemek, hedefi ve gerekli kapsamı
değiştirir; mevcut yanıtlayıcılara yeni katsayı vermek yeterli değildir.
\item Dosya, seçim ve yanıt süreçlerinden sonraki kayıttır. Ondan
alt örnek seçmek, seçilmeyen haneleri veya yanıt vermeyen kişileri
geri getirmez; ayrıca yeni bir seçim aşaması yaratır.
\item Ağırlıklı payda yetişkin alt nüfusun büyüklüğüne ilişkin tahmin;
ağırlıksız payda yetişkin kayıt sayısıdır. Nihai ağırlıktaki ek
düzeltmeler nedeniyle temel $1/\pi_i$ ile son ağırlık ayrılmalıdır.
\item Önce tam tasarım bilgisi korunur, sonra yetişkin alt evreni
tanımlanır. Varyans hesabı tabaka ve PSU katkılarıyla yapılır; oranı
ağırlıklandırmak gözlemleri bağımsız hale getirmez.
\item Örnek: ``Bu plan, belirtilen dönem ve hedef nüfustaki yetişkinler
için 60 yaş ve üzeri oranını tahmin etmeyi amaçlar. Sonuç, bugünkü
Türkiye nüfusunun oranı veya yaşın herhangi bir sonuç üzerindeki
nedensel etkisi olarak yorumlanamaz.'' Sayısal sonuç ayrıca doğrulanmalıdır.
\end{enumerate}

\textbf{Değerlendirme ölçütleri (10 puan).} Her görev iki ayrı
bileşenden değerlendirilir; bileşen başına doğru ve gerekçeli yanıt
1 puan, eksik veya yanlış yanıt 0 puandır.
\begin{enumerate}
\item Hedef--alt evren--dönem tanımı (1); kapsam değişikliğinin
ağırlık değişikliğinden ayrılması (1).
\item Çerçeve ile yanıtlayıcı dosyasının ayrılması (1); yeniden
seçimin kayıp kapsamı geri getirmediğinin açıklanması (1).
\item İki paydanın doğru yorumu (1); temel ve nihai ağırlık ayrımı (1).
\item Alt evren için tasarım bilgisinin korunması (1); kümeli ve
tabakalı varyans gereğinin açıklanması (1).
\item Kapsamı doğru belirten rapor cümlesi (1); desteklenmeyen
genelleme ve nedensellik iddialarının dışlanması (1).
\end{enumerate}
Doğru formülü yazıp yanlış evrene genellemek tam puan getirmez.

\newcommand{\VtwoN}{9254}
\newcommand{\VtwoAdults}{5569}
\newcommand{\VtwoOlder}{2150}
\newcommand{\VtwoDf}{15}
\newcommand{\VtwoWeighted}{29{,}1516}
\newcommand{\VtwoUnweighted}{38{,}6066}
\newcommand{\VtwoSE}{1{,}5529}
\newcommand{\VtwoLower}{25{,}8416}
\newcommand{\VtwoUpper}{32{,}4615}
\subsection{Sayısal uygulama: dosya oranı ve hedef nüfus oranı}
\label{sec:v2-nhanes-sonuclar}

Kaynakta \VtwoN{} kişi kaydı vardır; bunların \VtwoAdults{} tanesi
20 yaş ve üzerinde, \VtwoOlder{} tanesi 60 yaş ve üzerindedir.
Seçili alanlarda eksiksizlik, pozitif görüşme ağırlıkları ve 15
tabaka--30 PSU yapısı kontrol edilmiştir. Diğer sütunların eksikleri
nedeniyle kayıt silinmemiştir. Dosyanın sayısal kabulü, kamu dosyasının
gizli hane seçim çerçevesini içerdiği anlamına gelmez.

\begin{center}
\begin{tabular}{@{}lr@{}}
\toprule
Özet & Sonuç \\
\midrule
Ağırlıksız yetişkin kayıt oranı & \%\VtwoUnweighted \\
Ağırlıklı alt evren oranı & \%\VtwoWeighted \\
Taylor standart hatası & \VtwoSE{} yüzde puan \\
Tasarım serbestlik derecesi & \VtwoDf \\
Yaklaşık \%95 güven aralığı & [\%\VtwoLower; \%\VtwoUpper] \\
\bottomrule
\end{tabular}
\end{center}

Ağırlıklı pay yaklaşık 69,596 milyon, payda 238,738 milyondur;
bunlar dosyadaki kişi sayısı değil, görüşme ağırlıklarıyla elde edilen
nüfus toplamı tahminleridir. Ağırlıksız oran yanıtlayıcıların bileşimini,
ağırlıklı oran hedef alt nüfusu özetler. Farklı olmaları hesap çelişkisi
veya iki bağımsız araştırma kanıtı değildir.

\textbf{İzlenebilir varyans hesabı.} $B=\sum_iw_id_i$ ve
$z_i=w_id_i(y_i-\widehat p_w)/B$ olsun. Tabaka $h$ içindeki PSU $j$
için $Z_{hj}$ kişi katkılarının toplamıdır. Nihai PSU yaklaşımıyla,
sonlu evren düzeltmesi kullanmadan,
\[
\widehat{\Var}(\widehat p_w)=
\sum_h\frac{m_h}{m_h-1}\sum_j(Z_{hj}-\overline Z_h)^2
\]
hesaplanır. Bütün kayıtlar tasarımda kalır; çocukların alt evren
katkıları sıfırdır. Bu veride yetişkinler bütün 15 tabaka ve 30 PSU'da
bulunduğundan serbestlik derecesi $30-15=15$ alınır
\citep{nchsvariance}. PSU kodları tabaka koduyla birlikte kullanılır.

Aralık $\widehat p_w\pm t_{0{,}975;15}\widehat{\SE}$ biçiminde,
yaklaşık t-Wald öğretim aralığıdır. NCHS'nin bütün oran yayımlama
ölçütlerinin uygulandığı iddia edilmez. Kapsam veya yanıtsızlık
yanlılığı yalnız standart hata hesabıyla ölçülmüş olmaz.

\begin{researchbox}
\textbf{Örnek rapor.} NHANES 2017--2018 görüşme ağırlıkları ve
maskelenmiş varyans tasarımıyla, hedefteki kurum dışı sivil yetişkin
nüfusta 60 yaş ve üzeri oranı yaklaşık \%29,15 tahmin edilmiştir
(Taylor standart hatası 1,55 yüzde puan; yaklaşık \%95 t aralığı
\%25,84--\%32,46). Bu tahmin bugünkü nüfusa, Türkiye'ye veya kurum
nüfusuna doğrudan taşınamaz; nedensel etkiyi ölçmez. CDC'nin belirli
bir resmî tahmininin tekrarı olarak sunulmaz.
\end{researchbox}

\begin{softwarebox}
Tam duyarlıklı sonuçlar \path{companion/vakalar/v2-nhanes/results.json},
30 PSU'nun katkıları aynı klasörde \texttt{psu.csv} içindedir.
\texttt{export.py} çıktıları üretir; \texttt{verify\_delivery.py}
dosya bütünlüğünü, sayısal eşleşmeyi ve kitap bağlantısını denetler.
İkinci hesap pay/payda toplamlarının kovaryansından delta yöntemiyle
aynı varyansa ulaşmıştır. XPT okuyucusu ve t kritik değeri ortaktır;
R/SPSS veya bağımsız bir anket yazılımı çalıştırılmış sayılmaz.
Sayısal eşleşme, örnekleme kapsamının veya yanıt sürecinin kusursuzluğunu
kanıtlamaz.
\end{softwarebox}

\section{Çözümlü senaryo}

Bir araştırmacı üç sınıf düzeyindeki öğrencilerin okul aidiyeti puanlarını
karşılaştıracaktır. Her sınıf düzeyinden öğrenci seçmek istendiği için sınıf
düzeyi tabaka olarak kullanılabilir. Yalnız gönüllülerin katılması durumunda
her tabakada dahi gönüllülük yanlılığı kalabilir. Bu nedenle seçim ve yanıt
oranları ayrı ayrı raporlanmalıdır.

\section{Yazılım uygulaması: rastgele ve tabakalı seçim}

\begin{softwarebox}
\begin{lstlisting}
import pandas as pd

students = pd.DataFrame({
    "id": range(1, 13),
    "sinif": [1]*4 + [2]*4 + [3]*4
})

simple = students.sample(n=6, random_state=2026)
stratified = (students.groupby("sinif", group_keys=False)
              .sample(n=2, random_state=2026))

print(simple)
print(stratified.groupby("sinif").size())
\end{lstlisting}
İkinci seçim her sınıf düzeyinden iki öğrenciyi güvenceye alır. Gerçek bir
uygulamada seçim olasılıkları ve örnekleme çerçevesinin kaynağı da
belgelenmelidir.
\end{softwarebox}

\begin{assumptionbox}
Rastgele sayı üretmek ancak güncel ve hedef evreni kapsayan bir örnekleme
çerçevesi üzerinde anlamlıdır. Listede bulunmayan birimlerin seçilme
olasılığı sıfırdır ve bu kapsam sorunu yazılım tarafından giderilemez.
\end{assumptionbox}

\section{Çözümlü problemler}

\subsection*{Problem 1: Basit rastgele örnekleme ve seçilme olasılığı}

Bir fakültede 1.200 öğrenci bulunmaktadır. Güncel öğrenci listesinden basit
rastgele örneklemeyle 120 öğrenci seçilecektir. Her öğrencinin seçilme
olasılığını, temel örnekleme ağırlığını ve örnekleme oranını bulun.

\textbf{Çözüm.} Basit rastgele örneklemede bütün öğrencilerin dâhil edilme
olasılığı aynıdır:
\[
\pi_i=\frac{120}{1200}=0{,}10.
\]
Her öğrencinin örneğe seçilme olasılığı yüzde 10'dur. Temel örnekleme ağırlığı
\[
w_i=\frac{1}{0{,}10}=10
\]
olur. Örnekleme oranı da $n/N=0{,}10$'dur. Tasarım eksiksiz uygulanır ve
yanıtlamama olmazsa örneklemdeki her öğrenci evrende yaklaşık on öğrenciyi
temsil eder.

\textbf{Varsayım kontrolü.} Öğrenci listesinin güncel ve tekilleştirilmiş
olması gerekir. Listede bulunmayan öğrencilerin seçilme olasılığı sıfırdır;
yinelenen kayıtların ise seçilme olasılığı yapay biçimde yükselir.

\subsection*{Problem 2: Sistematik örnekleme}

Numaralandırılmış 2.400 öğrencilik bir listeden 120 öğrenci sistematik
örneklemeyle seçilecektir. Örnekleme aralığını bulun. Rastgele başlangıç 7 ise
ilk dört ve son seçilen sıra numarasını yazın.

\textbf{Çözüm.} Örnekleme aralığı
\[
k=\frac{N}{n}=\frac{2400}{120}=20
\]
olur. Rastgele başlangıç 1 ile 20 arasından seçilen 7 olduğuna göre sıra
numaraları 7, 27, 47, 67, $\ldots$ biçiminde ilerler. Son sıra numarası
\[
7+(120-1)20=2387
\]
olur.

Listenin sıralanışında her 20 birimde tekrarlanan bir örüntü varsa sistematik
seçim yanlı sonuç üretebilir. Örneğin liste her sınıfta aynı başarı sırasına
göre düzenlenmişse periyodiklik kontrol edilmelidir.

\subsection*{Problem 3: Orantılı tabakalı örnekleme}

Bir üniversitedeki 1.000 öğrencinin 600'ü lisans, 300'ü yüksek lisans ve
100'ü doktora öğrencisidir. Toplam 200 kişilik orantılı tabakalı örneklemde
her gruptan kaç kişi seçilmelidir?

\textbf{Çözüm.} Her tabakanın evrendeki payı örneklem hacmiyle çarpılır:
\[
n_{\text{lisans}}=200\frac{600}{1000}=120,
\]
\[
n_{\text{yüksek lisans}}=200\frac{300}{1000}=60,
\qquad
n_{\text{doktora}}=200\frac{100}{1000}=20.
\]
Böylece toplam $120+60+20=200$ kişi seçilir. Her tabaka aynı yüzde 20
oranında örneklendiğinden temel ağırlıklar eşittir. Her tabakanın içinde
öğrenciler rastgele seçilmelidir; tabaka kotasını dolduran ilk gönüllüleri
almak olasılıklı örnekleme değildir.

\subsection*{Problem 4: Küme örneklemesinde tasarım etkisi}

Yirmi sınıftan sınıf başına ortalama 25 öğrenci seçilerek toplam 500 öğrenciye
ulaşılmıştır. Küme içi korelasyonun yaklaşık $\rho=0{,}08$ olduğu varsayılsın.
Tasarım etkisini, yaklaşık etkin örneklem hacmini ve standart hata çarpanını
hesaplayın.

\textbf{Çözüm.} Ortalama küme büyüklüğü $m=25$ olduğundan
\[
\DEFF\approx1+(25-1)(0{,}08)=1+1{,}92=2{,}92.
\]
Yaklaşık etkin örneklem hacmi
\[
n_{\mathrm{etkin}}\approx\frac{500}{2{,}92}\approx171
\]
olur. Kümeleşmeye göre düzeltilmiş standart hata, basit rastgele örnekleme
standart hatasının yaklaşık
\[
\sqrt{2{,}92}\approx1{,}71
\]
katıdır. Dosyada 500 öğrenci bulunmasına rağmen belirsizlik, 500 bağımsız
öğrenci varmış gibi hesaplanmamalıdır.

\begin{mistakebox}
Kümeleşme gözlem sayısını silmez; bilgi miktarını ve standart hatayı etkiler.
Çözüm, sınıf içindeki öğrencilerden rastgele bazılarını veri setinden çıkarmak
değil, örnekleme tasarımına uygun analiz kullanmaktır.
\end{mistakebox}

\subsection*{Problem 5: Yanıtlamama ve düzeltme ağırlıkları}

İki fakülteden eşit sayıda, toplam 400 öğrenci seçilmiştir. A fakültesinde
200 öğrencinin 160'ı, B fakültesinde 200 öğrencinin 120'si ankete yanıt
vermiştir. Fakülte bazında yanıt oranlarını ve basit yanıtlamama düzeltme
ağırlıklarını bulun.

\textbf{Çözüm.} Yanıt oranları
\[
r_A=\frac{160}{200}=0{,}80,
\qquad
r_B=\frac{120}{200}=0{,}60
\]
olur. Yanıtlamama düzeltme çarpanları bu oranların tersidir:
\[
a_A=\frac{1}{0{,}80}=1{,}25,
\qquad
a_B=\frac{1}{0{,}60}\approx1{,}67.
\]
Bu ağırlıklarla A fakültesindeki 160 yanıt yaklaşık $160(1{,}25)=200$, B
fakültesindeki 120 yanıt ise yaklaşık $120(1{,}67)=200$ seçilmiş öğrenciyi
temsil eder.

Bu düzeltme, her fakülte içinde yanıtlayan ve yanıtlamayan öğrencilerin
araştırma değişkeni bakımından yeterince benzer olduğu varsayımına dayanır.
Yanıtlamama örneğin memnuniyetsizlikle doğrudan ilişkiliyse yalnız fakülte
ağırlığı bütün yanlılığı gideremeyebilir.

\begin{outputguidebox}
Örnekleme raporunda hedef evreni, çerçevenin kaynağını, seçim aşamalarını,
başlangıç ve son örneklem hacimlerini, yanıt oranını, ağırlıkları ve küme ya da
tabaka değişkenlerini kontrol edin. Yalnız ``rastgele seçildi'' ifadesi
tasarımın yeniden üretilebilmesi için yeterli değildir.
\end{outputguidebox}

\section{Literatür verisiyle yazılım uygulamaları}
\label{sec:spss-b06}

\textbf{Amaç ve veri.} \texttt{b06.csv} içinde 395 kayıt, yalnız bu
uygulama için sonlu bir örnekleme çerçevesidir \citep{cortez2008data}.
\texttt{srs40} geri koymadan basit rastgele seçilen 40 kaydı;
\texttt{strat40} GP'den 35, MS'den 5 kayıt seçimini gösterir. Seçimler sabit
tohumla önceden üretilmiştir.

\textbf{Ortak veri.} Doğrulanan B06 uygulamasında üç yazılım da
\nolinkurl{companion/spss/csv/b06.csv} dosyasını kullanmıştır. Dosya
395 satır ve 13 sütun içerir; B01'in öğrenci değişkenlerine seçim
göstergeleri ve \texttt{strw} ağırlığı eklenmiştir. B01 dosyasını yeniden
adlandırmak yeterli değildir. Seçimler yeniden üretilmez; aynı hazır
göstergeler kullanılır. Bu karşılaştırma Excel içe aktarmanın doğrulandığı
anlamına gelmez. Kodları \texttt{companion} klasörünü içeren paket kökünden
çalıştırın; veri sözlüğü ve hazırlık için bkz. sayfa~\pageref{sec:spss-guide}.

\subsection{IBM SPSS uygulaması}
\textbf{Başlamadan önce.} Aşağıdaki komut dosyası B06 CSV verisini
açar ve analiz eder. Menü yolunu kullanacaksanız önce veri açma ve
tanımlama komutlarını \texttt{EXECUTE} satırı dahil çalıştırın.
Dosya açma hatası alırsanız \texttt{/FILE} yolunu ve dosyanın gerçekten
bu konumda bulunduğunu kontrol edin; veri açılmadan sonraki komutlar
çalışmaz.

\begin{spssadimlar}
\spssadim{Henüz filtre ve ağırlık yokken \emph{Analyze $\to$ Descriptive Statistics $\to$ Descriptives} yolunda \texttt{g3} için ortalama ve standart sapmayı alın; tam dosyada 395 kayıt bulunduğunu kontrol edin.}
\spssadim{\emph{Data $\to$ Select Cases} yolunda \emph{If condition is satisfied $\to$ If} seçin; \texttt{srs40=1} yazın. \emph{Continue} ardından \emph{Filter out unselected cases} ve \emph{OK} seçin; kayıtları silmeyin.}
\spssadim{\emph{Analyze $\to$ Descriptive Statistics $\to$ Descriptives} yolunda \texttt{g3} için ortalama ve standart sapmayı alın; analizde 40 kayıt kullanıldığını kontrol edin.}
\spssadim{\emph{Data $\to$ Select Cases} penceresinde önce \emph{All cases} ile filtreyi kapatın. Ardından koşulu \texttt{strat40=1} yaparak yeniden filtreleyin.}
\spssadim{\emph{Analyze $\to$ Descriptive Statistics $\to$ Frequencies} yolunda \texttt{school} tablosunu alın. Henüz ağırlık vermeden GP=35, MS=5 sayımlarını kontrol edin.}
\spssadim{\emph{Data $\to$ Weight Cases} yolunda \emph{Weight cases by} seçin, \texttt{strw} değişkenini ağırlık alanına taşıyıp \emph{OK} seçin. \emph{Descriptives} ile \texttt{g3} ortalamasını tekrar alın.}
\spssadim{Bitirirken \emph{Data $\to$ Weight Cases $\to$ Do not weight cases} ve \emph{Data $\to$ Select Cases $\to$ All cases} seçin. Hem menü hem aşağıdaki komut yolu 395 kaydı korur; kayıtları kalıcı olarak silmeyin.}
\end{spssadimlar}

{\raggedright
\noindent\textbf{Komutların görevi.} \texttt{FILTER BY} seçili kayıtları analiz eder, diğerlerini silmez. \texttt{FILTER OFF} tüm kayıtları yeniden etkinleştirir. \texttt{WEIGHT BY} ağırlıklı nokta tahminini hesaplatır; işlem sonunda \texttt{WEIGHT OFF} ile kapatılır.\par}

\textbf{Uygulama kodu:} \texttt{b06}. Doğrulanan veri dosyası
\texttt{companion/spss} altındaki \texttt{csv/b06.csv} dosyasıdır.

\begin{lstlisting}[style=prismSPSS]
SET DECIMAL=DOT.
GET DATA /TYPE=TXT
 /FILE='companion/spss/csv/b06.csv'
 /ENCODING='UTF8'
 /ARRANGEMENT=DELIMITED /DELCASE=LINE
 /DELIMITERS="," /QUALIFIER='"'
 /FIRSTCASE=2
 /VARIABLES=id F8.0 school F8.0 sex F8.0 age F8.0
 studytime F8.0 absences F8.0 g1 F8.0 g2 F8.0
 g3 F8.0 pass10 F8.0 srs40 F8.0 strat40 F8.0
 strw F20.12.
FILTER OFF.
WEIGHT OFF.
SPLIT FILE OFF.
VARIABLE LABELS
 g3 'Yil sonu matematik notu'
 srs40 'Basit rastgele orneklem secimi'
 strat40 'Tabakali orneklem secimi'
 strw 'Tabaka agirligi'.
VALUE LABELS school 1 'GP' 2 'MS'.
VALUE LABELS srs40 strat40 0 'Secilmedi' 1 'Secildi'.
VARIABLE LEVEL
 id school sex pass10 srs40 strat40 (NOMINAL)
 studytime (ORDINAL)
 age absences g1 g2 g3 strw (SCALE).
FORMATS g3 strw (F12.6).
EXECUTE.
DESCRIPTIVES VARIABLES=g3 /STATISTICS=MEAN STDDEV.
FILTER BY srs40.
DESCRIPTIVES VARIABLES=g3 /STATISTICS=MEAN STDDEV.
FILTER OFF.
FILTER BY strat40.
FREQUENCIES VARIABLES=school.
WEIGHT BY strw.
DESCRIPTIVES VARIABLES=g3 /STATISTICS=MEAN.
WEIGHT OFF.
FILTER OFF.
EXECUTE.
\end{lstlisting}

\begin{outputguidebox}
\textbf{Paylaşılan SPSS çıktısıyla doğrulanan değerler.}
Basit rastgele seçilen 40 kaydın not ortalaması \SPSrsMean'dir.
Tabakalı seçimde GP için $349/35$, MS için $46/5$ ağırlıkları kullanıldığında
ortalama \SPStratMean\ bulunur. Karşılaştırma değeri, dosyanın tamamının
\SPMatMean\ ortalamasıdır. Tek bir çekimde hangisinin daha yakın olduğu,
yöntemlerden birinin her zaman daha iyi olduğunu kanıtlamaz.
Ağırlıklı çıktıda gösterilen toplam ağırlık 395'tir; gerçek seçilmiş
öğrenci sayısı yine 40'tır. \texttt{WEIGHT BY} burada yalnız ağırlıklı nokta
tahminini göstermek içindir, karmaşık örnekleme standart hatası üretmez.
\end{outputguidebox}


\par\noindent\begin{minipage}{\linewidth}
\subsection{R uygulaması}
\label{sec:r-gercek-b06}

\texttt{subset} seçili kayıtları ayrı nesneye alır; \texttt{weighted.mean} tabaka ağırlıklı ortalamayı verir.

\begin{lstlisting}[style=prismR]
veri <- read.csv("companion/spss/csv/b06.csv")
stopifnot(nrow(veri) == 395, ncol(veri) == 13,
          !anyNA(veri))
basit <- subset(veri, srs40 == 1)
tabakali <- subset(veri, strat40 == 1)
stopifnot(nrow(basit) == 40, nrow(tabakali) == 40)
ozet <- function(degerler)
  c(n=length(degerler), mean=mean(degerler),
    sd=sd(degerler))
print(cbind(tam_dosya=ozet(veri$g3),
            basit_rastgele=ozet(basit$g3)))
okul <- factor(tabakali$school, c(1, 2), c("GP", "MS"))
frekans <- table(okul)
print(frekans); print(100 * prop.table(frekans))
print(nrow(tabakali)); print(sum(tabakali$strw))
print(weighted.mean(tabakali$g3, tabakali$strw))
\end{lstlisting}

\textbf{Çıktıyı okuma.} Paylaşılan R çıktısı 395 satır, 13 sütun ve
her sütunda sıfır eksik değer gösterir. \texttt{n}, \texttt{mean} ve
\texttt{sd} satırları tam dosya ile basit rastgele örneklemi karşılaştırır.
Tabakalı örneklemde okul frekansları ağırlıksızdır; 40 kayıt ile 395
ağırlık toplamı ayrı yazdırılır. Ana \texttt{veri} nesnesi değişmez.
\end{minipage}\par

\par\noindent\begin{minipage}{\linewidth}
\subsection{Python uygulaması}
\label{sec:python-b06}

\texttt{loc} ve \texttt{copy} seçimi ayrı veri çerçevesinde tutar. \texttt{np.average} ağırlıklı ortalamayı hesaplar.

\begin{lstlisting}[style=prismPythonApp]
import pandas as pd
import numpy as np

veri = pd.read_csv("companion/spss/csv/b06.csv")
assert veri.shape == (395, 13)
assert not veri.isna().any().any()
basit = veri.loc[veri.srs40 == 1].copy()
tabakali = veri.loc[veri.strat40 == 1].copy()
assert len(basit) == len(tabakali) == 40
print(pd.DataFrame({
    "tam_dosya": veri.g3.agg(["count", "mean", "std"]),
    "basit_rastgele": basit.g3.agg(["count", "mean", "std"])
}))
okul = tabakali.school.map({1: "GP", 2: "MS"})
frekans = okul.value_counts().reindex(["GP", "MS"])
print(frekans); print(100 * frekans / len(tabakali))
print(len(tabakali)); print(tabakali.strw.sum())
print(np.average(tabakali.g3, weights=tabakali.strw))
\end{lstlisting}

\textbf{Çıktıyı okuma.} Paylaşılan Python çıktısı da 395 satır,
13 sütun ve her sütunda sıfır eksik değer gösterir. \texttt{count}, R'deki
\texttt{n}; \texttt{std} ise \texttt{sd} satırına karşılık gelir.
Okul frekansları, yüzdeler ve ağırlıklı ortalama R ve SPSS ile uyumludur.
Bu uygulama bir hipotez testi veya $p$ değeri üretmez; ana veri korunur.
\end{minipage}\par

\subsection{Sonuçların karşılaştırılması ve ortak rapor}
12 Eylül 2026'da paylaşılan SPSS tabloları ile R ve Python metin
çıktıları karşılaştırılmıştır. Tam dosya ve basit rastgele örneklemin
gözlem sayısı, ortalaması ve standart sapmasından oluşan altı hücre ile
tabakalı seçimin ağırlıklı ortalaması, gösterilen yuvarlama
hassasiyetlerinde uyumludur. Okul frekansları, yüzdeler, tabakalı
örneklem büyüklüğü ve ağırlık toplamı da eşleşmiştir. Proje CSV'sinden
yapılan hesaplar bu değerleri destekler. Aşağıdaki tablolar çıktıların
kitap düzenine uyarlanmış özetidir; ekran görüntüsü değildir.

\begin{table}[H]
\centering
\caption{B06'da tam dosya ve seçilmiş örneklemlerin not özetleri}
\label{tab:b06-dogrulanmis-betimsel}
\small
\begin{tabular}{@{}lrrr@{}}
\toprule
İncelenen veri & Kayıt sayısı & Ortalama & Std. sapma \\
\midrule
Tam dosya & 395 & 10,415190 & 4,581443 \\
Basit rastgele örneklem & 40 & 10,275000 & 5,565404 \\
Tabakalı örneklem (ağırlıklı) & 40 & 10,456347 & --- \\
\bottomrule
\end{tabular}
\par\smallskip
{\footnotesize Ondalıklı değerler altı basamağa yuvarlanmıştır.
İlk iki satırın standart sapmaları sırasıyla 394 ve 39 böleniyle
hesaplanmıştır; standart hata değildir. Son satırda yalnız ağırlıklı
ortalama raporlanmıştır; çizgi sıfır anlamına gelmez. SPSS'in bu
satır için gösterdiği $N=395$ ağırlık toplamıdır, gerçek kayıt sayısı
40'tır.\par}
\end{table}

\begin{table}[H]
\centering
\caption{B06 tabakalı örnekleminde ağırlıksız okul frekansları}
\label{tab:b06-dogrulanmis-tabaka}
\begin{tabular}{@{}lrr@{}}
\toprule
Okul & Frekans & Yüzde \\
\midrule
GP & 35 & 87,5 \\
MS & 5 & 12,5 \\
Toplam & 40 & 100,0 \\
\bottomrule
\end{tabular}
\par\smallskip
{\footnotesize Yüzdelerin paydası seçilen 40 kayıttır. Okul değişkeninde
eksik değer olmadığından SPSS'in \emph{Percent} ve \emph{Valid Percent}
sütunları eşittir. Ağırlıklar bu tablo alındıktan sonra etkinleştirilir.\par}
\end{table}

\Needspace{15\baselineskip}
\begin{outputguidebox}
Tablo~\ref{tab:b06-dogrulanmis-betimsel} içindeki tam dosya ortalaması,
yalnız bu 395 kayıt eğitim amaçlı sonlu evren kabul edildiğinde bilinen
evren ortalamasıdır. Basit rastgele örneklemin 10,275000 ve tabakalı
örneklemin 10,456347 ortalamaları bu değerin iki ayrı tahminidir.
Bu tek seçimde tabakalı tahmin daha yakındır; bu durum yöntemin her
örneklemde daha yakın sonuç vereceğini kanıtlamaz.

Tablo~\ref{tab:b06-dogrulanmis-tabaka} içinde GP için 35, MS için 5
kayıt vardır. Ağırlıklar sırasıyla $349/35$ ve $46/5$ olduğundan
toplamları $35(349/35)+5(46/5)=395$ olur. Bu işlem 355 yeni öğrenci
eklemez. Ağırlıklı nokta tahmini, tasarıma uygun standart hata veya
güven aralığı hesabı değildir.

Seçim göstergeleri öğretim amacıyla hazırlanmıştır; kaynak araştırmanın
özgün örnekleme tasarımını göstermez. Yazılımların hesaplama uyumu,
395 kaydın daha geniş bir öğrenci evrenini temsil ettiğini kanıtlamaz.
\end{outputguidebox}

\textbf{Raporlama örneği (doğrulanmış çıktılarla).}
``395 kayıtlık öğretim çerçevesinin not ortalaması 10,415190'dır.
Basit rastgele seçilmiş 40 kaydın ortalaması 10,275000, standart sapması
5,565404 bulunmuştur. GP'den 35 ve MS'den 5 kayıt içeren ayrı tabakalı
örneklemde, tabaka ağırlıklarıyla hesaplanan ortalama 10,456347'dir.
Ağırlıkların toplamı 395 olmakla birlikte gerçek örneklem büyüklüğü
40'tır. Karşılaştırma tek seçime aittir; yöntemlerin genel üstünlüğü
veya daha geniş bir evrene genellenebilirlik iddiası taşımamaktadır.''


\textbf{Görev.} $\bar x_{\rm tab}=(349/395)\bar x_{GP}+(46/395)\bar x_{MS}$
ile sonucu elle doğrulayın. Bu seçimlerin kaynak araştırmanın özgün
örnekleme tasarımı olmadığını raporlayın. Sonraki uygulamaya geçmeden
filtre ve ağırlığın kapalı olduğunu, etkin dosyanın 395 kaydı koruduğunu
kontrol edin.

\section{R ile çözümlü uygulama}
\label{sec:r-b06}

\textbf{Soru.} Her sınıfta dört öğrenci bulunan 12 kişilik çerçeveden altı
kişilik basit rastgele örneklem ve her sınıftan ikişer kişilik tabakalı
örneklem seçin. Seçilme olasılıklarını ve ağırlıkları bulun.

\begin{lstlisting}[style=prismR]
set.seed(2026)
ogrenciler <- data.frame(
  id = 1:12, sinif = factor(rep(1:3, each = 4))
)
basit <- ogrenciler[sample.int(nrow(ogrenciler), 6), ]
indisler <- split(seq_len(nrow(ogrenciler)),
                  ogrenciler$sinif)
secilen <- unlist(lapply(indisler,
                        function(indis) sample(indis, 2)),
                  use.names = FALSE)
tabakali <- ogrenciler[secilen, ]
tabakali$secim_olasiligi <- 2 / 4
tabakali$agirlik <- 1 / tabakali$secim_olasiligi
print(basit)
print(tabakali)
table(tabakali$sinif)
sum(tabakali$agirlik)
stopifnot(nrow(basit) == 6,
          anyDuplicated(basit$id) == 0,
          all(table(tabakali$sinif) == 2))
\end{lstlisting}

\begin{outputguidebox}
\textbf{Çözüm ve kontrol değerleri.} İki tasarımda da altı farklı öğrenci
seçilir. Tabakalı seçimde sınıf frekansları kesin olarak $(2,2,2)$'dir;
basit rastgele seçim bunu garanti etmez. Her öğrencinin seçilme olasılığı
$6/12=2/4=0{,}5$, tasarım ağırlığı 2'dir. Altı ağırlığın toplamı 12'dir.
\end{outputguidebox}

\textbf{Yorum.} Bu örnekte tabakalar eşit büyüklüktedir ve eşit sayıda
birim seçilir. Eşit olmayan tabaka hacimlerinde her tabaka için $N_h/n_h$
ayrı hesaplanmalıdır. Ağırlıklı ortalama hesaplamak, tasarıma uygun standart
hata hesabını kendiliğinden sağlamaz.

\textbf{Pekiştirme ve yanıt.} Bir sınıfın çerçeveden tamamen çıkarılması
kapsam hatası yaratır; kalan öğrencileri rastgele seçmek bu hatayı düzeltmez.

\section{Bölüm özeti}

\begin{itemize}
  \item Olasılıklı örneklemede birimlerin seçilme olasılıkları bilinir.
  \item Tabakalama temsili, küme örneklemesi veri toplama verimliliğini destekleyebilir.
  \item Olasılıksız örneklemede genelleme sınırları açıkça belirtilmelidir.
  \item Büyük örneklem rastgele hatayı azaltır; kapsam ve yanıtlamama yanlılığını gidermez.
  \item Eşit olmayan seçilme ve yanıt olasılıkları analizde ağırlıklandırma
  gerektirebilir.
  \item Küme içi benzerlik tasarım etkisini ve standart hatayı artırarak etkin
  örneklem hacmini azaltabilir.
\end{itemize}

\section*{Alıştırmalar}

\begin{enumerate}
  \item Hedef evren ile örnekleme çerçevesinin farklı olduğu bir örnek verin.
  \item Tabakalı ve küme örneklemesini amaçları bakımından karşılaştırın.
  \item Örneklem hacmini büyütmenin yanlılığı neden gidermediğini açıklayın.
  \item 5.000 birimlik bir evrenden 250 birim basit rastgele seçildiğinde
  seçilme olasılığını ve temel ağırlığı hesaplayın.
  \item $N=1800$ ve $n=90$ için sistematik örnekleme aralığını bulun; rastgele
  başlangıç 12 ise ilk beş sıra numarasını yazın.
  \item 700, 200 ve 100 kişilik üç tabakadan toplam 300 kişilik orantılı
  örneklem için tabaka hacimlerini belirleyin.
  \item Tabakalı örnekleme ile kota örneklemesi arasındaki temel farkı açıklayın.
  \item Ortalama küme büyüklüğü 30 ve küme içi korelasyon 0,05 ise yaklaşık
  tasarım etkisini hesaplayın.
  \item $n=600$ ve $\DEFF=2{,}5$ için yaklaşık etkin örneklem hacmini bulun.
  \item Bir ankette başlangıç örneklemi 800, yanıtlayan sayısı 520 ise genel
  yanıt oranını ve basit düzeltme çarpanını hesaplayın.
  \item Yanıtlamama ağırlığının yanıtlayanlarla yanıtlamayanlar sistematik
  olarak farklı olduğunda neden yetersiz kalabileceğini tartışın.
  \item Çok aşamalı bir ulusal okul araştırması için il, okul, sınıf ve öğrenci
  düzeylerini içeren bir seçim planı yazın.
\end{enumerate}